% GENERATED FILE. Edit canonical lessons, then run npm run generate:books.
% canonical-source-hash: fnv1a64:524f2f306a7c2cba
% canonical-lessons: GU-C52-kaka, GU-C52-kaki, GU-C52-mama, GU-C52-mami, GU-C52-husband

\chapter{Uncles, Aunts and a Husband}
\label{ch:gu-uncles-aunts-and-a-husband}
\hlchaptermodality{52}

\begin{chapteropening}
\textbf{By the end of this chapter:} \emph{I can name uncles and aunts on both sides, and a husband.}

Complete the last lesson of chapter 52: I can name uncles and aunts on both sides, and a husband.
\end{chapteropening}

\section[\texorpdfstring{\gu{કાકા} (kākā)}{કાકા (kākā)}]{\gu{કાકા} (kākā) — uncle (father's brother)}

\label{lesson:GU-C52-kaka}



\begin{quote}
Before the new one: say the Gujarati for grandmother, then the Gujarati for a child.
\end{quote}

\subsection*{You'll want to know: \gu{કાકા}}
\textbf{\gu{કાકા}} (\emph{kākā}) — \textquotedblleft{}uncle (father's brother)\textquotedblright{}.

\begin{cousinweb}
\textbf{\gu{કાકા}} (\emph{kākā}) goes back to \emph{*kākka}, \textquotedblleft{}senior male relative\textquotedblright{}, a form CDIAL reconstructs and takes from Dravidian, comparing Kannada \emph{kakka}, \textquotedblleft{}uncle\textquotedblright{} (CDIAL 2998).
\end{cousinweb}

\begin{grammarlens}[title={What you've built}]
One more word for uncles, aunts and a husband.
\end{grammarlens}

\subsection*{Your turn}
\begin{itemize}
\raggedright
  \item \emph{Say it:} \emph{kākā}
  \item \emph{Say it:} \emph{kākā}, once more
  \item \emph{Say it:} say \emph{kākā}
  \item \emph{Recall:} say the Gujarati for grandmother, then the Gujarati for a child, then say \emph{kākā} again
\end{itemize}

\begin{tcolorbox}[breakable,colback=teal!4,colframe=teal!35!black,title={Before you move on}]
What does \gu{કાકા} mean? (\textquotedblleft{}Uncle (father's brother)\textquotedblright{}.) Say it once more.
\end{tcolorbox}
\section[\texorpdfstring{\gu{કાકી} (kākī)}{કાકી (kākī)}]{\gu{કાકી} (kākī) — aunt (father's brother's wife)}

\label{lesson:GU-C52-kaki}



\begin{quote}
Before the new one: say the Gujarati for a child, then the Gujarati for uncle.
\end{quote}

\subsection*{You'll want to know: \gu{કાકી}}
\textbf{\gu{કાકી}} (\emph{kākī}) — \textquotedblleft{}aunt (father's brother's wife)\textquotedblright{}.

\begin{cousinweb}
\textbf{\gu{કાકી}} (\emph{kākī}) is the feminine of \emph{kākā}, \textquotedblleft{}uncle\textquotedblright{}. CDIAL lists both under \emph{*kākka}, \textquotedblleft{}senior male relative\textquotedblright{}, which it takes from Dravidian (CDIAL 2998).
\end{cousinweb}

\begin{grammarlens}[title={What you've built}]
One more word for uncles, aunts and a husband.
\end{grammarlens}

\subsection*{Your turn}
\begin{itemize}
\raggedright
  \item \emph{Say it:} \emph{kākī}
  \item \emph{Say it:} \emph{kākī}, once more
  \item \emph{Say it:} say \emph{kākī}
  \item \emph{Recall:} say the Gujarati for a child, then the Gujarati for uncle, then say \emph{kākī} again
\end{itemize}

\begin{tcolorbox}[breakable,colback=teal!4,colframe=teal!35!black,title={Before you move on}]
What does \gu{કાકી} mean? (\textquotedblleft{}Aunt (father's brother's wife)\textquotedblright{}.) Say it once more.
\end{tcolorbox}
\section[\texorpdfstring{\gu{મામા} (māmā)}{મામા (māmā)}]{\gu{મામા} (māmā) — uncle (mother's brother)}

\label{lesson:GU-C52-mama}



\begin{quote}
Before the new one: say the Gujarati for uncle, then the Gujarati for aunt.
\end{quote}

\subsection*{You'll want to know: \gu{મામા}}
\textbf{\gu{મામા}} (\emph{māmā}) — \textquotedblleft{}uncle (mother's brother)\textquotedblright{}. Gujarati names the side of the family.

\begin{cousinweb}
\textbf{\gu{મામા}} (\emph{māmā}) is a tadbhava of Sanskrit \emph{māma}, \textquotedblleft{}uncle\textquotedblright{} (CDIAL 10055). CDIAL calls it a nursery word and notes similar words in Dravidian.
\end{cousinweb}

\begin{grammarlens}[title={What you've built}]
One more word for uncles, aunts and a husband.
\end{grammarlens}

\subsection*{Your turn}
\begin{itemize}
\raggedright
  \item \emph{Say it:} \emph{māmā}
  \item \emph{Say it:} \emph{māmā}, once more
  \item \emph{Say it:} say \emph{māmā}
  \item \emph{Recall:} say the Gujarati for uncle, then the Gujarati for aunt, then say \emph{māmā} again
\end{itemize}

\begin{tcolorbox}[breakable,colback=teal!4,colframe=teal!35!black,title={Before you move on}]
What does \gu{મામા} mean? (\textquotedblleft{}Uncle (mother's brother)\textquotedblright{}.) Say it once more.
\end{tcolorbox}
\section[\texorpdfstring{\gu{મામી} (māmī)}{મામી (māmī)}]{\gu{મામી} (māmī) — aunt (mother's brother's wife)}

\label{lesson:GU-C52-mami}



\begin{quote}
Before the new one: say the Gujarati for aunt, then the Gujarati for uncle.
\end{quote}

\subsection*{You'll want to know: \gu{મામી}}
\textbf{\gu{મામી}} (\emph{māmī}) — \textquotedblleft{}aunt (mother's brother's wife)\textquotedblright{}.

\begin{cousinweb}
\textbf{\gu{મામી}} (\emph{māmī}) is the feminine of \emph{māmā}. CDIAL lists both under Sanskrit \emph{māma}, \textquotedblleft{}uncle\textquotedblright{}, a nursery word with similar words in Dravidian (CDIAL 10055).
\end{cousinweb}

\begin{grammarlens}[title={What you've built}]
One more word for uncles, aunts and a husband.
\end{grammarlens}

\subsection*{Your turn}
\begin{itemize}
\raggedright
  \item \emph{Say it:} \emph{māmī}
  \item \emph{Say it:} \emph{māmī}, once more
  \item \emph{Say it:} say \emph{māmī}
  \item \emph{Recall:} say the Gujarati for aunt, then the Gujarati for uncle, then say \emph{māmī} again
\end{itemize}

\begin{tcolorbox}[breakable,colback=teal!4,colframe=teal!35!black,title={Before you move on}]
What does \gu{મામી} mean? (\textquotedblleft{}Aunt (mother's brother's wife)\textquotedblright{}.) Say it once more.
\end{tcolorbox}
\section[\texorpdfstring{\gu{પતિ} (pati)}{પતિ (pati)}]{\gu{પતિ} (pati) — a husband}

\label{lesson:GU-C52-husband}



\begin{quote}
Before the new one: say the Gujarati for uncle, then the Gujarati for aunt.
\end{quote}

\subsection*{You'll want to know: \gu{પતિ}}
\textbf{\gu{પતિ}} (\emph{pati}) — \textquotedblleft{}a husband\textquotedblright{}.

\begin{cousinweb}
\textbf{\gu{પતિ}} (\emph{pati}) is Sanskrit \emph{pati}, \textquotedblleft{}master, husband\textquotedblright{}, taken over whole as a learned loan, a tatsama (CDIAL 7727).
\end{cousinweb}

\begin{grammarlens}[title={What you've built}]
One more word for uncles, aunts and a husband.
\end{grammarlens}

\subsection*{Your turn}
\begin{itemize}
\raggedright
  \item \emph{Say it:} \emph{pati}
  \item \emph{Say it:} \emph{pati}, once more
  \item \emph{Say it:} say \emph{pati}
  \item \emph{Recall:} say the Gujarati for uncle, then the Gujarati for aunt, then say \emph{pati} again
\end{itemize}

\begin{tcolorbox}[breakable,colback=teal!4,colframe=teal!35!black,title={Before you move on}]
What does \gu{પતિ} mean? (\textquotedblleft{}A husband\textquotedblright{}.) Say it once more.
\end{tcolorbox}
